% TOP_PROBLEM: {"id": 3231, "title": "Partial List Coloring", "category_id": 3, "category": {"id": 3, "name": "graph_theory", "display_name": "Graph Theory", "description": "Problems involving graphs, networks, and their properties.", "slug": "graph-theory", "order_index": 3, "created_at": "2026-07-31T15:26:25.671Z"}, "status": "open", "problem_number": "OPG-771", "set_id": 12}
% FIRST_POSTED: 2026-10-01
% ATTEMPT_STATUS: unresolved
% MODEL: GPT 6 Astra Ultra
% UnsolvedMath ID 3231; source catalogue code OPG-771
% !TeX program = lualatex
% Research attempt dated 2026-10-01; canonical statement preserved.
\documentclass[11pt,a4paper]{article}
\usepackage[margin=25mm]{geometry}
\usepackage{fontspec}
\setmainfont{lmroman10-regular.otf}[BoldFont=lmroman10-bold.otf,
 ItalicFont=lmroman10-italic.otf,BoldItalicFont=lmroman10-bolditalic.otf]
\usepackage{amsmath,amssymb,mathtools}
\usepackage{unicode-math}
\setmathfont{latinmodern-math.otf}
\AtBeginDocument{\let\setminus\smallsetminus}
\usepackage{xurl,hyperref}
\hypersetup{colorlinks=true,urlcolor=blue}
\setcounter{secnumdepth}{0}
\setlength{\parindent}{0pt}
\setlength{\parskip}{5pt}
\setlength{\emergencystretch}{3em}
\begin{document}
\section*{Partial List Coloring}\label{3231}
{\small UnsolvedMath ID 3231 · OPG-771 · Graph Theory · OpenGarden}
\subsection{Definitions and mathematical statement}
Let $G=(V,E)$ be a nonempty
finite simple graph with
$n=|V|$. A list assignment
gives each vertex $v$ a
finite set $L(v)$ of allowed
colours. An $L$-colouring
is a function $c$ with
$c(v)\in L(v)$ at every
vertex and $c(u)\neq c(v)$
on every edge. The list
chromatic number $\chi_\ell(G)$
is the least integer $q$
such that every list assignment
with at least $q$ colours
at each vertex permits a
proper $L$-colouring.

Put $q=\chi_\ell(G)$, so
$q\geq1$, and let $t$
be an integer with
$0\leq t\leq q$. A partial
proper list colouring consists
of a subset $U\subseteq V$
and a proper list colouring
of the induced graph $G[U]$
using the original lists.
Vertices outside $U$ are
left uncoloured, so they
impose no colour choice.

The conjecture asserts that
for every assignment of
exactly $t$ colours to each
vertex there is such a
partial colouring with
\[
 |U|\geq\frac{tn}{q},
 \quad\text{equivalently}\quad
 |U|\geq\left\lceil\frac{tn}{q}\right\rceil.
\]
The lists may differ arbitrarily,
and $U$ may be selected
after seeing them. The cases
$t=0$ and $t=q$ respectively
allow the empty colouring
and require a full colouring
already guaranteed by the
definition. Intermediate list
sizes request the stated
linear proportion uniformly
over all graphs and all
lists, not merely in
expectation.
\subsection{Short English statement}
With lists of t colours, can every graph colour at least the fraction t divided by its list chromatic number of its vertices?
\subsection{Research attempt}
\paragraph{Scope and method.}
This research attempt by GPT-6 Astra Ultra is dated 1 October 2026.
It preserves the catalogue statement and distinguishes elementary proofs
below from cited prior work. No novelty claim or formal proof-assistant
verification is made. The final paragraph states the precise remaining scope.
\paragraph{Current literature changes the historical target.}
Noel's September 2026 preprint [R1] reports a fourteen-vertex graph with
$\chi_\ell=3$ and a two-list assignment colouring at most nine vertices.
Since $9<28/3$, this disproves the historical assertion. We preserve the
canonical statement as a catalogue record; it must not be interpreted as
a verified current open-status claim. This attempt develops independently
proved lower bounds and checks the finite bad-list calculation, without
claiming the counterexample as a new discovery.

\paragraph{Endpoints and singleton lists.}
The cases $t=0$ and $t=q$ follow respectively from the empty colouring
and the definition of $q$-choosability. For $t=1$, an ordinary proper
$q$-colouring exists by assigning the same $q$ colours to every vertex.
One colour class has size at least $\lceil n/q\rceil$ and is independent.
All its vertices can use their respective singleton list colours, regardless
of whether those colours agree. Thus the requested lower bound is always
valid for singleton lists. The selected class need not coincide with
any one of the colours appearing in those lists.

\paragraph{An exact class where the proposed ratio holds.}
Suppose $G$ is a disjoint union of cliques of orders $m_1,\ldots,m_s$.
For arbitrary lists of size $t$, select $\min(t,m_i)$ vertices in clique
$i$. Their lists have a system of distinct representatives: any nonempty
subfamily has union of size at least $t$, at least the number of lists
in that subfamily. Hall's theorem applies. This colours
$\sum_i\min(t,m_i)$ vertices. Giving all vertices the same $t$ colours
shows this guarantee is exact, because no clique can use more than $t$
vertices. Here $q=\max_i m_i$; the lower bound follows from
\[
            \min(t,m_i)\ge\frac{tm_i}{q}\qquad(0\le t\le q).
\]
Indeed if $m_i\le t$, use $t\le q$; otherwise use $m_i\le q$.
Equal-sized cliques attain the ratio, showing that the target could never
have been universally increased even within this positive family.

\paragraph{A universal degree-based lower bound.}
For $t\ge1$, independently choose at each vertex one of its $t$ list
colours uniformly and assign independent continuous random priorities.
Retain a vertex if its priority is smaller than that of every neighbour
which chose the same colour. Retained adjacent vertices cannot have equal
colours: of any such pair, the one with larger priority is discarded.
Thus the retained set is always properly list-coloured.

Conditional on the colour chosen at $v$, let $X_v$ count neighbours
choosing it. Each contributes probability at most $1/t$, so
$\mathbb E X_v\le d(v)/t$. Conditional on these choices, retention has
probability $1/(1+X_v)$. Convexity of $x\mapsto1/(1+x)$ gives
\[
 \mathbb P(v\text{ retained})\ge\frac{t}{t+d(v)}.
\]
By summing expectations, some deterministic outcome colours at least
$\left\lceil\sum_v t/(t+d(v))\right\rceil$ vertices. This is an actual
existence bound for arbitrary lists. Its denominator involves degrees,
so it does not imply the historical $tn/q$ claim in general.

\paragraph{Verification and the failed padding route.}
The root batch script implements the graph and bad lists explicitly
specified in [R1] and exhaustively computes a maximum partial colouring
of size nine. This checks that finite list assignment, not universal
three-choosability; the latter is supplied by the paper's proof.
Separately, padding every list by $q-t$ fresh common colours allows a
full colouring, but deleting vertices using fresh colours has no automatic
proportional guarantee. A fresh colour can be used on a large independent
set, so counting the number of fresh colours does not bound the number
of vertices lost.

\paragraph{Final status.}
\textbf{UNRESOLVED as a new-proof attempt.} The historical affirmative
statement is reported false in [R1], with its finite bad-list bound
rechecked here. The singleton, clique-union and probabilistic lower bounds
above are proved. No new resolution, optimal replacement constant, or
general characterization of counterexamples is claimed.

\subsection{Sources}
\begin{itemize}
\item[S1] ULAM.AI and the UnsolvedMath Contributors, supplied problems.json, record 3231 (OPG-771), OpenGarden collection. Catalogue identity and source transcription; CC BY 4.0.
\url{https://huggingface.co/datasets/ulamai/UnsolvedMath}.
\item[S2] Open Problem Garden, Partial List Coloring. Original problem and discussion; the mathematical formulation below is expanded from the supplied source record.
\url{http://www.openproblemgarden.org/op/partial_list_coloring}.
\item[R1] Jonathan A. Noel, \emph{The Partial List Colouring
Conjecture is False}, arXiv:2609.23291v1, 20 September 2026.
Primary full text checked 1 October 2026. The result is prior work.
\url{https://arxiv.org/html/2609.23291v1}.
\end{itemize}
\end{document}
